%LaTeX-inställningar%%%%%%%%%%%%%%%%%%%%%%%%%%%%%%%%%%%%%%%%%%%%%%%%%%%
% Kompliera med pdflatex!!!
\documentclass[a4paper]{dtek}

\setcounter{secnumdepth}{5}

\titleformat{\paragraph}[hang]{\normalfont\normalsize\bfseries}{\theparagraph}{1em}{}
\titlespacing*{\paragraph}{0pt}{3.25ex plus 1ex minus 0.2ex}{0.7em}

%FYLL I VID ÄNDRINGAR%%%%%%%%%%%%%%%%%%%%%%%%%%%%%%%%%%%%%%%%%%%%%%%%%%
\newcommand{\updated}{2024-11-22} %Insert

%%%%%%%%%%%%%%%%%%%%%%%%%%%%%%%%
% PRISBASBELOPP
% sätt till nuvarande prisbasbelopp för att annotera med prisbasbelopp i spalten till höger.
% OBS detta ska vara avstängt för det faktiska reglementesdokumentet och när ändringar ska godkännas på sektionsmöten
%\newcommand{\nuvarandeprisbasbelopp}{47600}


\newcommand{\prisbasbelopp}[1]{
    #1 prisbasbelopp
    \ifdef{\nuvarandeprisbasbelopp}{
        \marginpar{
            \small{ \textbf{\roundandprint{\xintiexpr#1 * \nuvarandeprisbasbelopp\relax}kr}}
        }
    }{}
}


\title{Ekonomiskt reglemente}
\date{Uppdaterad \updated}
%Dokumentstart%%%%%%%%%%%%%%%%%%%%%%%%%%%%%%%%%%%%%%%%%%%%%%%%%%%%%%%%%
\begin{document}

\makeheadfoot%

\begin{center}
  {\Huge \textbf{Ekonomiskt Reglemente för Datateknologsektionen}}

  \includegraphics[width=300pt]{dteklogo.pdf}

  {\LARGE Chalmers, Göteborg}

  \vspace{2cm}

  {\LARGE Uppdaterad \updated}
\end{center}

\newpage
\setcounter{tocdepth}{1}
\tableofcontents
\newpage

%%%%%%%%%%%%%%%%%%%%%%%%%%%%%%
%Allmänt
%%%%%%%%%%%%%%%%%%%%%%%%%%%%%%
\section{Allmänt}
\subsection{Definitioner}
Kommittéer används som ett samlingsnamn för kommittéer och sektionstyrelsen.

\subsection{Filosofi}
Detta reglemente har skapats för att skapa tydlighet över vad sektionens pengar får och inte får användas till. 
Det är därmed tänkt att vara ett hjälpmedel för sektionstyrelsen, revisorerna och sektionens kommittéer i deras arbete. 
Grundfilosofin är att sektionens pengar är till för sektionens samtliga medlemmar, vilket i praktiken innebär att arrangemang som inte är öppna för dessa inte bör belasta sektionens ekonomi. 
Viktigt att tänka på vid tolkandet utav reglementet är att det inte har utformats för att i onödan begränsa kommittéers möjligheter eller lust till roliga arrangemang och upptåg.
\subsection{Prisbasbelopp}
\begin{itemize}
  \item Det prisbasbelopp som ska användas under hela verksamhetsåret är det prisbasbelopp som är aktuellt vid verksamhetsårets början.
  \item Delar av prisbasbelopp ska avrundas till närmaste hundratal.
\end{itemize}
%%%%%%%%%%%%%%%%%%%%%%%%%%%%%%%%
%Kommittéer
%%%%%%%%%%%%%%%%%%%%%%%%%%%%%%%%
\section{Kommittéer}
\label{sec:kommitteer}

\subsection{Ordförandes skyldigheter}
Ordförande i kommitté är skyldig att
\begin{itemize}
\item Kontinuerligt meddela kommitténs ekonomiska status till styret
\item Tillsammans med kommitténs kassör ansvara för att kommittén förvaltar sina tillgångar i enlighet med sektionens stadgar, reglementen och beslut.
\item Tillsammans med kommitténs kassör skriva ett bindande avtal med styret angående skuldfrågan vid felaktig bokföring
\end{itemize}
\subsection{Kassörs skyldigheter}
\label{sec:kommittee_kassor}
Kassör i kommitté är skyldig att
\begin{itemize}
\item Föra kassabok av sådan typ som godkänts av sektionens revisorer
\item Tillsammans med kommitténs ordförande ansvara för att kommittén förvaltar sina tillgångar i enlighet med sektionens stadgar, reglementen och beslut.
\item Tillsammans med kommitténs ordförande skriva ett bindandeavtal med Styret angående skuldfrågan vid felaktigt förd bokföring
\item På varje sektionsmöte redovisa kommitténs ekonomiska situation
\item Arkivera kommitténs bokföring, på en plats anvisad av Styret, så lång tid som föreskrivs för den organisationsform som datateknologsektionen är.
\item Lägga budget enligt \secref{sec:budget}
\end{itemize}
\subsection{Avsaknande av kassör}
För kommittéer som saknar kassör ansvarar istället sektionskassören för samtliga av punkterna i \secref{sec:kommittee_kassor} med undantag för att lägga budget vilket görs i samråd med kommitténs ordförande.
\subsection{Budget}
\label{sec:budget}
Varje kommitté ska i början av året lägga en budget för hur dess pengar ska spenderas. Det är upp till sektionskassören att delge en mall för budget för att underlätta kommittéernas arbete. I budgeten ska det tydligt framgå hur pengarna ska spenderas med en uppskattning för vilken månad utgifter och inkomster hamnar i.

\subsubsection{Budgetförslag}
Varje kommittés budgetförslag ska godkännas av sektionstyrelsen, som rapporterar vidare relevant information till sektionsmötet.


\subsection{Kapital}
\label{sec:sektionsforeningar_startkapital}
\subsubsection{Startkapital}
Kommittéer ska vid mandatperiodens början inte ha några tillgångar.
\subsubsection{Startlån}
Följande Kommittéer har rätt att låna \prisbasbelopp{0.25} från sektionstyrelsen i form av startlån:
\begin{itemize}
  \item DRust
  \item DAG
  \item Delta
  \item D6
  \item DNollK
  \item DBus
  \item DKock
\end{itemize}
DAG har rätt låna ytterligare \prisbasbelopp{1} i form av utökat startlån.
\subsubsection{Slutkapital}
Samtliga kommitteers tillgångar och skulder tillfaller sektionstyrelsen vid mandatperiodens slut.
\subsubsection{Stödkapital}
Styret har rätt att utföra kortfristiga lån till kommittéer vid behov. Dessa ska isåfall betalas tillbaka senast i slutet av mandatperioden.

%%%%%%%%%%%%%%%%%%%%%%%%%%%%%
% Intresseföreningar
%%%%%%%%%%%%%%%%%%%%%%%%%%%%%
\section{Intresseföreningar}
\subsection{Ordförande}
Ordförande i intresseförening är skyldig att
\begin{itemize}
  \item Kontinuerligt meddela intresseförenings ekonomiska status till styret.
  \item Tillsammans med intresseförenings kassör ansvara för att intresseförening förvaltar sina tillgångar i enlighet med sektionens stadgar, reglementen och beslut.
  \item Tillsammans med intresseförenings kassör skriva ett bindande avtal med styret angående skuldfrågan vid felaktig bokföring.
\end{itemize}

\subsection{Kassör}\label{IntresseKassör}
Kassör i intresseförening är skyldig att
\begin{itemize}
  \item Tillsammans med intresseförenings ordförande ansvara för att intresseförening förvaltar sina tillgångar i enlighet med sektionens stadgar, reglementen och beslut.
  \item Tillsammans med intresseförenings ordförande skriva ett bindandeavtal med Styret angående skuldfrågan vid felaktigt förd bokföring.
  \item På varje sektionsmöte redovisa intresseförenings ekonomiska situation.
  \item Arkivera intresseförenings bokföring, på en plats anvisad av sektionsstyret, så lång tid som föreskrivs för den organisationsform som datateknologsektionen är.
  \item Lägga budget enligt \ref{IntresseBudget}
\end{itemize}

\subsection{Avsaknad av kassör}
För intresseföreningar som saknar kassör ansvarar istället sektionskassören för samtliga av punkterna i \ref{IntresseKassör} med undantag för att lägga budget vilket görs i samråd med intresseförenings ordförande

\subsection{Budget}\label{IntresseBudget}
Varje intresseförening ska i början av året lägga en budget för hur dess pengar ska spenderas.
Det är upp till sektionskassören att delge en mall för budget för att underlätta intresseföreningarnas arbete.
I budgeten ska det tydligt framgå hur pengarna ska spenderas med en uppskattning för vilken månad utgifter och inkomster hamnar i.

\subsubsection{Rapportering}
Varje intresseförenings budgetförslag ska godkännas av sektionstyrelsen, som rapporterar vidare relevant information till sektionsmötet

\subsection{Kapital}
\subsubsection{Startkapital}
Intresseförening skall vid mandatperiodens början inte ha några tillgångar.
\subsubsection{Slutkapital}
Samtliga intresseföreningars tillgångar och skulder tillfaller sektionstyrelsen vid mandatperiodens slut.
\subsubsection{Stödkapital}
Styret har rätt att utföra kortfristiga lån till intresseföreningar vid behov.
Dessa ska isåfall betalas tillbaka senast i slutet av mandatperioden.

%%%%%%%%%%%%%%%%%%%%%%%%%%%%%
% Äskning
%%%%%%%%%%%%%%%%%%%%%%%%%%%%%
\section{Äskning}
\subsection{Till sektionstyrelsen}
Medlem, kommittéer eller intresseföreningar som önskar extra ekonomiska medel till sin verksamhet eller arrangemang ska inkomma med önskemål och skäl till styret.
Styret kan bevilja extra medel om beloppet faller inom ramarna för de olika fonderna som sektionstyrelsen förvaltar. Om äskningen inte fallar inom dessa ramar kan den ändå beviljas förutsatt att beloppet understiger \prisbasbelopp{0.25} och inte överstiger sektionens budget.
\subsection{Till sektionsmötet}
Medlem, kommittéer eller intresseföreningar kan också inkomma med önskemål om extra ekonomiska medel till sin verksamhet eller arrangemang genom motion till sektionsmötet. Finansieringen av motionen ska göras med prioritet på medel från de fonder som sektionstyrelsen förvaltar.

%%%%%%%%%%%%%%%%%%%%%%%%%%%%%%
% Förmåner
%%%%%%%%%%%%%%%%%%%%%%%%%%%%%%
\section{Rättigheter}
För rättigheter får inte pengar spenderas på alkohol om inte annat anges.

\subsection{Internrepresentation}
\label{sec:internreps}
Varje kommitté har inte rätt att belasta sin ekonomi med mer än:
\begin{itemize}
    \item[-] \prisbasbelopp{0.004} för teambuildingaktiviteter per person.
    \item[-] \prisbasbelopp{0.0020} (avrundat till närmaste femkrona) för överlämning per person som deltar från avgående samt påstigande år.
    \item[-] \prisbasbelopp{0.005}  per post för aspning och aspplagg sammanlagt.
\end{itemize}

\subsubsection{Alkohol inom teambuilding}
Upp till en tredjedel av beloppet för teambuildingaktiviteter får användas till alkohol per person i kommittén.

\subsubsection{Använding av teambuildingbudget}
För att få använda hela det budgeterade beloppet för teambuilding måste 50\% av beloppet användas inom 8 månader från mandatperiodens början. Annars får inte mer än det som spenderats alternativt upp till 25\% användas.

\subsubsection{Ökad aspningsbudget}
Varje kommitée har rätt att skicka in en förfrågan för att använda upp till
0.017 prisbasbelopp per person i kommitéen för aspning och aspplagg sammanlagt. För att
få lov att spendera denna summa på aspning ska en kommitté få godkännande av
sektionstyrelsen innan inköp. Detta genom att kommittén skickar in plan för aspningen samt
en kostnadskalkyl som sektionstyrelsen baserar sitt beslut på.


\subsection{Representationskläder}
Varje kommitté har inte rätt att belasta sin ekonomi mer än \prisbasbelopp{0.01} per person i kommitén
för att synliggöra sig själva på sektionen och campus. 
Representationskläder kan ses som overaller, arbetsbyxor, hoodies, t-shirts eller liknande.

\subsubsection{Förfrågan om ökad representationsklädsbudget}
Varje kommitté har rätt att skicka in en förfrågan för att höja gränsen
per person i kommitteen till \prisbasbelopp{0.022}.
För att få lov att köpa in representationskläder för denna summa ska en 
kommitté få godkännande av sektionstyrelsen innan inköp. Detta genom att kommittén 
skickar förslag på kläder och en uppskattad kostnad som sektionstyrelsen baserar 
sitt beslut på.

\subsubsection{Ansökan}
En ansökan om representationskläder ska innehålla:
\begin{itemize}
  \item Det som man vill köpa in och vad som tidigare har köpts in.
  \item När och i vilka sammanhang representationskläder kommer användas.
  \item Varför gruppen vill ha representationskläder.
  \item En ungefärlig kostnadsberäkning.
\end{itemize}

\subsubsection{Godkännande}
När sektionstyrelsen beslutar huruvida ett förslag godkänns ska följande tas hänsyn till:
\begin{itemize}
    \item Att priset inte är onödigt dyrt (jämfört med liknande kläder för andra kommittéer).
    \item Att kläderna har en anknytning till verksamheten.
    \item Att kommittén kommer använda kläderna tillräckligt mycket för att motivera inköp.
\end{itemize}

\subsection{Representation \& Möteskostnader}
Vid arrangemang för sektionen kan representation eller möteskostnader beviljas. Dessa belastar kommitténs ekonomi. Hur mycket som godkänns baseras på arrangemangets längd där beloppet avrundas till närmaste femkrona:
\begin{itemize}
    \item 1 till 4 timmar: \prisbasbelopp{0.00075} per person.
    \item 4 till 8 timmar: \prisbasbelopp{0.0015} per person.
    \item Mer än 8 timmar: \prisbasbelopp{0.00225} per person.
\end{itemize}
\subsubsection{Redovisning}
Vid redovisning av representation eller möteskostnader ska deltagare samt tiden de deltagit finnas med för att belopp ska kunna godkännas.

\subsection{Skyldigheter}
Ovanstående belopp får inte överskrida kommitténs ekonomiska resurser. \\
%%%%%%%%%%%%%%%%%%%%%%%%%%%%%%%
% Sponsring
%%%%%%%%%%%%%%%%%%%%%%%%%%%%%%%
\section{Sponsring}
\subsection{Samråd}
Om kommittéer eller intresseföreningar vill söka sponsring ska detta göras i samråd med datas arbetsmarknadsgrupp, DAG.
\section{Fonder}
De fonder som beskrivs i denna del syftar till sparande av medel för framtida användande enligt respektive fonds syfte, antingen planerat eller spontant. Ett uttag ur en fond ska inte belasta Sektionstyrelsens godkända budget, däremot ska inköp och investeringar vara föremål för redovisning i balans– respektive resultaträkning. Detta gäller om inget annat specificeras i fonden.

\subsection{Uttag ur fonder}
Sektionsmötet äger alltid rätt att göra uttag ur fonder i enhet med fondens syfte. Sektionsmötet kan genom en kvalificerad majoritet göra uttag ur fonder i strid mot fondens syfte, ett välmotiverat behov till uttaget ska protokollföras vid ett sådant beslut. Alla andra uttagsformer måste specificeras i det ekonomiska reglementet.

\subsection{Redovisning av uttag}
Alla uttag ur fonder ska redovisas på sektionsmöte i sektionstyrelsens verksamhetsrapport.

\subsection{Kapitalfonden}
\subsubsection{Syfte}
\label{sec:kapitalfond_syfte}
Syftet med kapitalfonden är att avlasta sektionens, kommittéernas och intresseföreningars ekonomi från stora investeringar. Pengarna ska användas till saker som har ett bestående värde, samt är till gagn för sektionens medlemmar direkt eller indirekt. Pengar ska ej användas till driftbidrag eller stöd för förgänglig verksamhet. Fonden ska ej användas till verksamhet som lokalfond är ämnad för.
\subsubsection{Förvaltning}
Fonden förvaltas av sektionens styrelse.
\subsubsection{Avsättning}
\begin{itemize}
\item En av sektionstyrelsen budgeterad summa som godkänts av sektionsmötet
\item All avkastning ifrån kapitalfonden under verksamhetsåret.
\end{itemize}
\subsubsection{Uttag}
\begin{itemize}
\item Sektionstyrelsen har rätt att bevilja uttag ur fonden på belopp upp till totalt \prisbasbelopp{0.25} per tillfälle. Uttag av belopp överstigande \prisbasbelopp{0.25} ska godkännas av sektionsmöte innan medel utbetalas.
\item Sektionstyrelsen äger inte rätt att ta ut mer än 50\% av fondens totala värde per tillfälle.
\end{itemize}

\subsection{Lokalfonden}
\subsubsection{Syfte}
\label{sec:lokalfond_syfte}
Syftet med lokalfonden är att säkra medel för underhåll och reparationer av sektionens lokaler, samt för inköp av inventarier. Pengarna ska användas till större reparationer och ommålningar av sektionslokalerna, samt möbler till trivselytor där teknologen i gemen har tillträde.
\subsubsection{Förvaltning}
Fonden förvaltas av sektionens styrelse.
\subsubsection{Avsättning}
\label{sec:lokalfond_avsattning}
\begin{itemize}
\item Minst Tio (10)\% av under verksamhetsåret influtna sektionsavgifter tillförs lokalfonden.
\item En av sektionstyrelsen budgeterad summa som godkänts av sektionsmötet
\item All avkastning ifrån lokalfonden under verksamhetsåret.
\end{itemize}
\subsubsection{Uttag}
\begin{itemize}
\item Sektionstyrelsen disponerar sjuttiofem (75) \% av årets tillförda kapital,
enligt \secref{sec:lokalfond_avsattning}, för basdrift av sektionslokalerna.
\item Sektionstyrelsen har rätt att besluta om ytterliggare uttag, dock ska detta redovisas inför nästkommande sektionsmöte. Vid detta sektionsmöte ska i sådana fall även genomförda eller planerade inköp, reparationer och underhåll redovisas.
\end{itemize}

\subsection{Bilfonden}
\subsubsection{Syfte}
\label{sec:bilfond_syfte}
Syftet med bilfonden är att bygga upp en buffert för bilinköp och minska belastningen av sektionens ekonomi vid omfattande skador.
\subsubsection{Förvaltning}
Fonden förvaltas av sektionens styrelse.
\subsubsection{Avsättning}
\begin{itemize}
\item En av sektionstyrelsens budgeterad summa som godkänts av sektionsmötet.
\item All avkastning ifrån bilfonden under verksamhetsåret.
\end{itemize}
\subsubsection{Uttag}
\begin{itemize}
    \item Sektionstyrelsen äger rätt att bevilja uttag för omfattande reparation av sektionensbil.
    \item Sektionstyrelsen har rätt att köpa in en ny bil efter att ha fått en proposition till sektionsmötet godkänd med en specifikation på kostnad och krav för bilen.
\end{itemize}

\subsection{Idéfonden}
\subsubsection{Syfte}
\label{sec:idefond_syfte}
Syftet med idéfonden är att möjliggöra för datateknologer som önskar medel till arrangemang eller inventarier vars effekt är till gagn för hela teknologsektionens medlemmar.
\subsubsection{Förvaltning}
Idéfonden förvaltas av sektionstyrelsen.
\subsubsection{Avsättning}
\begin{itemize}
\item Tio (10) \% av under verksamhetsåret influtna sektionsavgifter tillförs idéfonden.
\item En av sektionstyrelsen budgeterad summa som godkänts av sektionsmötet.
\item All avkastning ifrån idéfonden under verksamhetsåret.
\end{itemize}
\subsubsection{Uttag}
\begin{itemize}
  \item Sektionstyrelsen har rätt att, på förslag av medlem i teknologsektionen, bevilja uttag på upp till 0,25 prisbasbelopp som går i hand med idéfondens syfte.
  \item Sektionsmötet har rätt att, på förslag av medlem i teknologsektionen, bevilja uttagsom går i hand med \secref{sec:idefond_syfte}.
\end{itemize}

\subsection{Husfonden}
\subsubsection{Syfte}
\label{sec:husfond_syfte}
Syftet med husfonden är att möjliggöra framtida fastighetsinvestering eller lokalför-värvning för sektionensmedlemmars bästa.
\subsubsection{Förvaltning}
Husfonden förvaltas av sektionstyrelsen.
\subsubsection{Avsättning}
\begin{itemize}
  \item En av sektionstyrelsen budgeterad summa som godkänts av sektionsmötet.
  \item All avkastning ifrån husfonden under verksamhetsåret.
\end{itemize}
\subsubsection{Uttag}
\begin{itemize}
  \item Sektionstyrelsen har rätt att göra uttag som uppfyller syftet efter att ha fått en proposition till sektionsmötet godkänd som inkluderar en kostnadskalkyl och annan relevant information om investeringen.
\end{itemize}
\subsection{Jubileumsfonden}
\subsubsection{Syfte}
\label{sec:jubileumfond_syfte}
Syftet med Jubileumsfonden är att möjliggöra firande av sektionens framtida jubileum, genom att möjliggöra större och dyrare evenemang.
\subsubsection{Förvaltning}
Jubileumsfonden förvaltas av sektionstyrelsen.
\subsubsection{Avsättning}
\begin{itemize}
  \item En av sektionstyrelsen budgeterad summa som godkänts av sektionsmötet.
  \item All avkastning ifrån Jubileumsfonden under verksamhetsåret.
\end{itemize}
\subsubsection{Uttag}
\begin{itemize}
  \item Sektionstyrelsen har rätt att göra uttag under jubileumsår.
\end{itemize}
\end{document}
